\documentclass[10pt]{article}
\usepackage[margin=20mm]{geometry}
\usepackage{amsmath,amssymb,mathtools,fontspec,unicode-math,xurl,hyperref}
\IfFontExistsTF{DejaVu Serif}{\setmainfont{DejaVu Serif}}{\setmainfont{Latin Modern Roman}}
\setmathfont{latinmodern-math.otf}
\setlength{\parindent}{0pt}\setlength{\parskip}{6pt}\sloppy
\title{Centering and variance signs in the maximum
of an i.i.d. characteristic polynomial}
\author{}
\date{}
\begin{document}
\maketitle
Abstract: The retained June 2024 preprint of Cipolloni and Landon is examined. For the real Ginibre ensemble, an exact determinant second moment and a holomorphic submean estimate contradict the mesoscopic-strip conclusion of Theorem 2.3 with its printed centering. The probability of one of its asserted lower-bound events is at most \(C n^{-1/5}\), although the theorem requires it to tend to one. A separate variance formula in Lemma 6.2 has a negative leading term for an allowed characteristic. Both signs remain in the selectively compared retained February 2026 preprint. Internal formulas suggest bounded corrections, but a corrected maximum theorem is not proved here. The findings are restricted to the acquired versions and the specified claims.\par
\section*{Identity, version, and scope}
The target is M25-0408, Maximum of the Characteristic Polynomial of I.I.D. Matrices, by Giorgio Cipolloni and Benjamin Landon, DOI `10.1002/cpa.22250`. The primary source is the 87-page arXiv:2405.05045v2, submitted 6 June 2024 and bearing the title-page date 7 June 2024 [CL2]. Printed and PDF page numbers agree. All unqualified source locators refer to this file. The retained 88-page arXiv:2405.05045v3 was submitted 17 February 2026 and bears the title-page date 18 February 2026 [CL3]. Only its identity and the two relevant statements were compared; it was not reviewed in full. Neither file is represented as a visually verified journal version of record.\par
The full primary source, including Appendices A–F and the bibliography, was read. The central chain runs from Hermitization and local laws to Gaussian characteristic increments, branching-random-walk comparisons, and the removal of regularization and Gaussian components. Detailed checks below concern the real-case normalization, an exact moment obstruction, and the deterministic variance input. Reading the other arguments is not represented as an independent reconstruction of every local law, coupling, comparison estimate, or cited theorem. The retained acquisition records describe bounded supplement and code searches; a missing identified repository is neither a proof of global absence nor a mathematical defect. No author code was executed.\par
Write \(p_n(z)=\det(X-zI)\) and \(P_n(z)=\log|p_n(z)|\). Definition 2.1 [p. 11] permits independent real Gaussian entries \(X_{ij}=n^{-1/2}g_{ij}\), \(g_{ij}\sim N(0,1)\): they are centered, have variance \(1/n\), and their rescaled moments of every fixed order are uniformly finite. This ensemble supplies the concrete witness. In the bulk, source (2.3) gives \[\varphi(z)=\frac{|z|^2-1}{2}.\]\par
\section*{A main strip-maximum statement fails with the printed center}
Equation reference: eq:center, eq:claim. Problem 1.  Major. In source (2.5), p. 12, the center is \[
 E_n(z)=n\varphi(z)-\frac14\log q_n(z),
 \qquad q_n(z)=4(\operatorname{Im}z)^2+n^{-1}.\] Theorem 2.3(2.7) then asserts that, for every \(0<\alpha<1/2\), \(0<r<1\), and \(\varepsilon>0\), \[
 \mathbb P\left\{
 \left(\frac1{\sqrt2}-\varepsilon\right)\log n
 \le \max_{z\in S_n}\bigl(P_n(z)-E_n(z)\bigr)
 \le \left(\frac1{\sqrt2}+\varepsilon\right)\log n
 \right\}\longrightarrow1,\] where \[S_n=\{z\in\mathbb C: |z|\le r,\ n^{-\alpha}\le\operatorname{Im}z
                         \le2n^{-\alpha}\}.\] For real Ginibre matrices, \(\alpha=1/4\), \(r=1/2\), and \(\varepsilon=1/\sqrt2-3/5>0\), the probability in [eq:claim] (\#eq:claim) instead tends to zero.\par
\subsection*{An exact determinant second moment}
Equation reference: eq:moment. For independent centered entries with \(\mathbb E|X_{ij}|^2=1/n\), the principal-minor expansion is \[p_n(z)=\sum_{S\subseteq[n]}(-z)^{n-|S|}\det X[S,S],
 \qquad \det X[\varnothing,\varnothing]=1.\] Each permutation monomial in a minor contains distinct matrix entries. Two such monomials from different subsets or different permutations have an ordered entry appearing in only one of them. Independence and centering make their cross expectation zero, also when the second monomial is conjugated. When both the subset and permutation coincide, the expectation of the squared modulus is \(n^{-|S|}\). Thus \[\begin{aligned}
 \mathbb E|p_n(z)|^2
 &=\sum_{k=0}^n {n\choose k}k!\,n^{-k}|z|^{2(n-k)}
 \notag\\
 &=n!\,n^{-n}\sum_{j=0}^n\frac{(n|z|^2)^j}{j!}
 \le C\sqrt n\,e^{\,n(|z|^2-1)}.                 \end{aligned}\] The exponential-series bound gives a sum at most \(e^{n|z|^2}\). Stirling’s upper bound is \[n!\le C\sqrt n(n/e)^n.\] This identity uses only independence, centering, and the second moment; in particular it applies to the allowed real Gaussian witness. It is not an approximation to a logarithmic covariance.\par
\subsection*{A continuum maximum bound}
Put \[F_n(z)=e^{-n(|z|^2-1)}|p_n(z)|^2.\] Equation [eq:moment] (\#eq:moment) gives \(\mathbb EF_n(z)\le C\sqrt n\) uniformly in \(z\). A bound at one point would not control the maximum, so the following deterministic step is essential.\par
Equation reference: eq:submean. Fix \(z\) and let \(\rho=n^{-1/2}\). As a function of \(w\), \[g_z(w)=p_n(w)\exp\{-n\overline z(w-z)\}\] is entire. The submean inequality for \(|g_z|^2\), together with \[e^{-n(|z|^2-1)}|g_z(w)|^2
       =F_n(w)e^{n|w-z|^2},\] gives \[
 F_n(z)\le\frac{e}{\pi\rho^2}
                  \int_{B_\rho(z)}F_n(w)\,dA(w)
       =\frac{en}{\pi}\int_{B_\rho(z)}F_n(w)\,dA(w).\] Here the submean inequality follows directly by expanding the holomorphic function in a power series about \(z\): circle integration annuls distinct powers, and all remaining squared terms are nonnegative.\par
Equation reference: eq:supmoment. Let \(S_n^\rho\) be the Euclidean \(\rho\)-neighborhood of \(S_n\). Every ball in [eq:submean] (\#eq:submean) with \(z\in S_n\) is contained in this one deterministic set, so \[\sup_{S_n}F_n\le\frac{en}{\pi}\int_{S_n^\rho}F_n(w)\,dA(w).\] The neighborhood is contained in the rectangle with horizontal width \(2r+2\rho\) and vertical height \(n^{-\alpha}+2\rho\). Since \(\alpha<1/2\), \(|S_n^\rho|\le C_r n^{-\alpha}\). Tonelli’s theorem now gives \[
 \mathbb E\sup_{S_n}F_n\le C_r n^{3/2-\alpha}.\] No random net, derivative estimate, or unproved continuity modulus is used. The polynomial and \(F_n\) are continuous; zeros of the determinant cause no difficulty because the event below can be expressed using the continuous nonnegative quantity \(F_nq_n^{1/2}\).\par
\subsection*{A fully specified contradiction}
Equation reference: eq:tail. For \(z\in S_n\) and \(n\ge1\), \[q_n(z)\le16n^{-2\alpha}+n^{-1}\le17n^{-2\alpha}.\] The printed center [eq:center] (\#eq:center) implies \[e^{2(P_n(z)-E_n(z))}=F_n(z)q_n(z)^{1/2}.\] Consequently, for any fixed \(a\), \[\begin{aligned}
 &\mathbb P\left\{\max_{S_n}(P_n-E_n)\ge a\log n\right\}
 \notag\\
 &\quad\le
 \mathbb P\left\{\sup_{S_n}F_n
                 \ge17^{-1/2}n^{2a+\alpha}\right\}
 \le C_r n^{3/2-2\alpha-2a},                     \end{aligned}\] by Markov’s inequality and [eq:supmoment] (\#eq:supmoment). Take \(\alpha=1/4\), \(a=3/5\), and \(r=1/2\). The set \(S_n\) is nonempty for all sufficiently large \(n\), and \[\frac32-2\alpha-2a=-\frac15,\qquad
 \varepsilon=\frac1{\sqrt2}-\frac35>0.\] The two-sided event in [eq:claim] (\#eq:claim) is a subset of the event in [eq:tail] (\#eq:tail); hence its probability is at most \(C n^{-1/5}\to0\). All distributional, moment, spatial, and parameter hypotheses of the stated theorem are satisfied. A fixed-point mean heuristic alone would not prove this contradiction; [eq:submean] (\#eq:submean) is what closes the gap between moments and the continuum maximum.\par
\subsection*{Counterevidence, repair scope, and impact}
The sign was checked on the actual source image, not inferred from text extraction. It is repeated in the plus logarithmic correction of \(\Psi_n\) in (2.20)–(2.21), p. 14. There is substantial internal evidence of an intended opposite sign: the real-axis expression (2.8), p. 12, uses \(P_n-n\varphi+\frac14\log n\); the mean formula (4.3), p. 22, and Appendix E’s mean computation (E.41)–(E.44), p. 79, have the positive logarithmic correction in the expectation. Moreover, the passage (10.27)–(10.28), p. 53, adds \(\alpha\log n\) to the uncorrected doubled logarithm on this strip. That addition is compatible with the negative correction in \(\Psi_n\), whereas the printed (2.21) has a positive one.\par
Equation reference: eq:correctcenter. The natural candidate center is \[
 E_n^{\rm corr}(z)=n\varphi(z)+\frac14\log q_n(z).\] The following algebraic consequence is fully verified: \[E_n-E_n^{\rm corr}=-\frac12\log q_n
                   =\alpha\log n+O(1)
       \quad\hbox{uniformly on }S_n,\] because \(q_n/n^{-2\alpha}\) stays between positive fixed constants. Thus \[P_n-E_n^{\rm corr}
       =(P_n-E_n)+\alpha\log n+O(1).\] This is a leading-order change on the strip. If a corrected theorem with center [eq:correctcenter] (\#eq:correctcenter) and coefficient \(1/\sqrt2\) were established, it would give coefficient \(1/\sqrt2-\alpha\) with the old center. That statement is a conditional implication, not a theorem proved in this report.\par
The verified repair is the centering algebra and its agreement with the cited internal mean and real-axis formulas. A full repair requires consistent signs in the definitions and all real-case comparison, increment, and branching-random-walk steps. That reconstruction has not been completed here, and no author adoption is asserted. The same (2.5) and (2.7) remain visible on p. 12 of the retained later preprint [CL3].\par
The severity is Major because [eq:tail] (\#eq:tail) directly falsifies an explicit principal conclusion, source Theorem 2.3(2.7), for an allowed ensemble. The lower-bound argument in §11.1 [pp. 56–57] expressly concludes (2.7), using the chain through Theorems 9.5 and 10.7 [pp. 48, 53]. This finding does not itself refute the complex-case Theorem 2.2, the full-disk real maximum (2.6), or every method in the paper. In particular, a false stated strip result used to prove (2.6) does not logically establish that (2.6) is false. No claim about an uninspected journal-final formula is made.\par
\section*{The integrated variance has the reversed leading logarithm}
Problem 2.  Major. In Lemma 6.2(6.5), p. 28, the first logarithm in the integral of a deterministic Brownian variance is printed as \(\log(\eta_s/\eta_{S_1})\). For a permitted characteristic decreasing from a mesoscopic scale to a microscopic scale away from the real axis, the asserted expression is strictly negative of order \(\log n\). An integrated nonnegative variance cannot have this value.\par
\subsection*{The statement and all parameter choices}
Section 6 takes the real flow \(X_t=\sqrt{1-T}\,Y+n^{-1/2}B_t\), where \(Y\) is real i.i.d. and \(B_t\) has independent standard real Brownian entries. One can take \(Y\) real Ginibre. The self-consistent variance is \(1-T+t\), and is one at \(t=T\). The characteristic satisfies \[\partial_t\eta_t=-\operatorname{Im}m_t^z(i\eta_t),
 \qquad m_t^z(i\eta_t)\ \hbox{constant along it}\] [source (5.6) and the computation on pp. 29–30]. Fix \(z=i/2\), in any bulk disk of radius greater than \(1/2\). Choose an arbitrarily small fixed \(c\in(0,1/4)\) meeting the lemma’s sufficient-smallness conditions, and set \[\varepsilon_1=\varepsilon_2=c,\quad T=n^{-c},\quad
 \eta_T=n^{-1+c},\quad S_1=0,\quad s=S_2=T.\] The auxiliary \(\xi>0\) in the lemma can be chosen sufficiently small independently. The required \(\eta_T\ge n^{-1+\varepsilon_2}\) holds with equality, and \(0\le S_1<S_2\le T\).\par
At the endpoint of variance one, write \(m_T^z(i\eta_T)=iv_n\). The bulk cubic equation, or source (6.15), yields \[v_n=\sqrt{1-|z|^2}+O(\eta_T)
     =\frac{\sqrt3}{2}+O(\eta_T).\] For completeness, the cubic at \(\eta=0\) and \(m=iv\) reduces to \(1/v=v+|z|^2/v\), giving \(v^2=1-|z|^2\). The positive bulk root is simple, so the implicit function theorem gives the displayed \(O(\eta_T)\) expansion. Constancy along the characteristic implies exactly \[\eta_0=\eta_T+Tv_n\asymp n^{-c}.\] This characteristic stays positive and in the small bulk regime throughout \([0,T]\).\par
\subsection*{A nonnegative quantity with a negative asserted value}
Equation reference: eq:badvar. With \(d=|z-\overline z|^2=1\) and \(b=2\sqrt{1-|z|^2}=\sqrt3\), the printed (6.5) reads \[
 \int_0^T V_u\,du
 =\log\frac{\eta_T}{\eta_0}
  +\log\frac{1+b\eta_0}{1+b\eta_T}
  +O\left(T\log n+\frac{\log n}{n\eta_T}\right).\] The first term is \(-(1-2c)\log n+O(1)\). The second is \(O(n^{-c})\), and the remainder is \(O(n^{-c}\log n)=o(1)\). Therefore the asserted right side is negative for all sufficiently large \(n\).\par
On the left, \(V_u\) is deterministic and nonnegative: it occurs under the real square root in the Brownian representation (6.4). The proof moreover defines it in (6.18), p. 30, and explicitly states \(V_u\asymp1/\eta_u>0\). Thus the contradiction is already deterministic. No choice of coupling, high-probability exceptional event, or hidden constant in the stated remainder can repair a negative leading logarithm of diverging magnitude.\par
\subsection*{The bounded correction and its limitations}
The first logarithm should have the reversed orientation, \(\log(\eta_{S_1}/\eta_s)\). This is checked directly for the leading variance integral: \[\int_{S_1}^s
      \frac{\operatorname{Im}m_u^z(i\eta_u)}{\eta_u}\,du
   =-\int_{\eta_{S_1}}^{\eta_s}\frac{\,d\eta}{\eta}
   =\log\frac{\eta_{S_1}}{\eta_s}.\] Source (5.25), p. 26, already prints this orientation for the corresponding complex-characteristic variance. The additional real contribution is treated by (6.16), p. 30. Reversing the first logarithm removes the exhibited sign obstruction and agrees with these local integrations; it is not a claim to have re-proved every coupling and remainder in Lemma 6.2.\par
The same reversed ratio is retained in (6.5) on p. 29 of [CL3]. Conversely, the later use of positive variance scales in (6.33) and (6.45) [pp. 33, 35] is evidence that the intended probabilistic argument does not use a negative variance. This counterevidence limits the conclusion: the printed lemma is false, but the report does not infer failure of every subsequent upper bound. Lemma 6.2 is invoked explicitly for the real-case Gaussian tail estimates and the two increment variances used in §6.4. Its invalid displayed guarantee is therefore a core proof input, supporting Major severity under the claim-level standard, despite the local sign correction being simple.\par
\section*{Coverage, calculation, and repair boundaries}
The review distinguishes a disproved principal statement from an invalid core intermediate formula. The following scope is exact.\par
Finding | Established result | Verified repair boundary\par
Problem 1 | Source (2.7) is false with (2.5); a specified allowed event has probability at most \(C n^{-1/5}\). | Opposite-sign centering changes the strip value by \(\alpha\log n+O(1)\). The corrected maximum theorem and all its dependencies are not independently proved.\par
Problem 2 | Source (6.5) gives a negative integrated variance for an allowed characteristic. | The first logarithm must be reversed by direct characteristic integration. Full coupling and error estimates remain beyond this bounded repair.\par
Independent exact symbolic diagnostics check the determinant second-moment identity for dimensions \(1\) through \(4\), the weighted holomorphic gauge identity, the exponent \(-1/5\), positivity of the chosen \(\varepsilon\), and the opposite leading variance signs. They are retained with their source and actual execution output under this paper’s verification directory. The general determinant and continuum arguments are proved above; finite checks are supporting diagnostics, not substitutes for those proofs. The variance contradiction is analytic and does not depend on numerical simulation.\par
The full acquired source has been read to the end, with focused visual checks of the identity, hypotheses, exact signs, proof dependencies, internal counterevidence, and later-version comparison. Appendices A–D and F were read for their role in the comparison chain; their external results have not all been independently re-proved. No issue count was imposed. No assertion of correctness for every unreported displayed formula is implied by the absence of an additional finding.\par
\section*{Scope and references}
The primary PDF has 936,922 bytes and SHA256 \[\begin{gathered}
\texttt{d591169236ab70305bf1dc214b5dfbf53}\\
\texttt{d60a57289da7b43ce26214c76637fad}.
\end{gathered}\] Its source manifest has SHA256 \[\begin{gathered}
\texttt{e4c5c5f8c8302d25b3560dabc61c916a69}\\
\texttt{c9045ab26aa8d5394aefb4ec3e75f4}.
\end{gathered}\] The exact-current English TeX, genuine compiler output, retained source/compiler/dependency payloads, parse and render receipts, and individual actual page inspections are bound in the companion provenance and inventory files. Rendering alone is not recorded as viewing. This is bounded initial recovery work, not independent scientific acceptance. The scientific decision remains pending and there are no human reviews.\par
2 G. Cipolloni and B. Landon, Maximum of the Characteristic Polynomial of I.I.D. Matrices, arXiv:2405.05045v2, submitted 6 June 2024, retained 87-page PDF dated 7 June 2024. Target DOI: `10.1002/cpa.22250`.\par
G. Cipolloni and B. Landon, Maximum of the Characteristic Polynomial of I.I.D. Matrices, arXiv:2405.05045v3, submitted 17 February 2026, retained 88-page PDF dated 18 February 2026. Selective comparison of identity, (2.5)/(2.7), and Lemma 6.2(6.5) only.\par
\end{document}
